\section*{Chapitre \ref{ch:b3:electrons-in-solids} --- Les électrons dans les solides}

\begin{solution}{exo:b3:electrons-in-solids:1}
(a) $\tau = m_{\text{e}}/ne^2\rho = \qty{2.5e-14}{s}$. (b) $\ell =
v_{\text{F}}\tau \approx \qty{39}{nm}$. (c) Environ 110 \omterm{def:b3:crystalline-solids:lattice}{paramètres
de maille}. (d) Un électron qui file devant une centaine d'ions sans
diffuser ne peut pas rebondir sur les ions eux-mêmes --- le \omterm{def:b3:crystalline-solids:lattice}{réseau}
doit lui être transparent, exactement ce que le \omterm{thm:b3:electrons-in-solids:bloch}{théorème de Bloch}
(\cref{thm:b3:electrons-in-solids:bloch}) démontrera plus tard.
\end{solution}

\begin{solution}{exo:b3:electrons-in-solids:2}
(a) $v_{\text{d}} = I/neA = \qty{7.3e-4}{m/s}$ --- moins d'un
millimètre par seconde. (b) Environ 23 minutes. (c) Le \emph{champ}
s'établit le long du fil à une vitesse proche de celle de la lumière
et met toute la mer d'électrons en dérive d'un coup~: les marcheurs
sont lents, l'ordre de marcher ne l'est pas. (d) $v_{\text{d}}/v_{\text{F}} \sim
10^{-9}$~: la conduction est un biais imperceptible sur un mouvement
quantique violent, non un écoulement classique paisible.
\end{solution}

\begin{solution}{exo:b3:electrons-in-solids:3}
(a) Bande à moitié pleine~: \omterm{def:b3:electrons-in-solids:classes}{métal}, $\dd\rho/\dd T > 0$. (b) Deux
électrons rempliraient exactement sa bande --- le magnésium conduit
parce que sa bande pleine \emph{recouvre} la suivante, vide~: \omterm{def:b3:electrons-in-solids:classes}{métal},
$\dd\rho/\dd T > 0$. (c) Bande pleine, gap de \qty{5.5}{eV}~:
\omterm{def:b3:electrons-in-solids:classes}{isolant} (formellement $\dd\rho/\dd T < 0$, mais sans pratiquement
aucun porteur à compter). (d) \omterm{def:b3:electrons-in-solids:classes}{Semi-conducteur}~: $\dd\rho/\dd T < 0$,
le changement de signe exponentiel qui trahit la classe.
\end{solution}

\begin{solution}{exo:b3:electrons-in-solids:4}
(a) $\lambda = hc/E_{\text{g}} = \qty{1.1}{\micro m}$~: le silicium
est transparent dans l'infrarouge au-delà --- et c'est pourquoi les
caméras infrarouges peuvent avoir des lentilles en silicium. (b)
\qty{225}{nm}, dans l'ultraviolet~: aucun photon visible ne peut
être absorbé, donc le diamant est limpide comme l'eau. (c)
\qty{365}{nm}~: un gap assez large pour le bleu (\qty{2.8}{eV})
existe dans les nitrures et presque nulle part ailleurs de façon
pratique --- la DEL bleue a attendu des décennies la croissance des
matériaux. (d) Son gap fixe son photon~: la surcharge ajoute de la
chaleur, non de la largeur de gap, et la diode cuit dans le noir ---
la couleur d'une DEL est une constante du matériau, non un réglage
de luminosité.
\end{solution}

\begin{solution}{exo:b3:electrons-in-solids:5}
(a) $\kappa/\sigma T = 400/(5.9\times10^7\times300) =
\qty{2.3e-8}{W.\ohm/K^2}$. (b) À dix pour cent près de $L$. (c)
Les mêmes électrons de la surface de Fermi portent la charge et la
chaleur, si bien que leur temps de diffusion se simplifie dans le
rapport. (d) La casserole d'acier conduit la chaleur vers les
aliments parce que ses électrons bougent~; le manche de bois,
\omterm{def:b3:electrons-in-solids:classes}{isolant} dans les deux sens, les garde --- et garde la chaleur ---
loin de votre main.
\end{solution}

\begin{solution}{exo:b3:electrons-in-solids:6}
(a) Cinquante pour cent de plus --- flagramment absents de
l'expérience. (b) $T/T_{\text{F}} \approx 0.004$~: le terme
électronique est étranglé sous le pour cent. (c) Le $T^3$ du \omterm{def:b3:crystalline-solids:lattice}{réseau}
s'effondre plus vite que le $T$ des électrons~: sous quelques
kelvins, les électrons «~négligeables~» sont tout ce qui reste. (d) $C/T =
\gamma + \beta T^2$~: l'ordonnée à l'origine $\gamma$ pèse les
électrons, la pente $\beta$ le \omterm{def:b3:crystalline-solids:lattice}{réseau} --- un seul graphique, les
deux locataires.
\end{solution}

\begin{solution}{exo:b3:electrons-in-solids:7}
(a) $k_{\text{F}} = (3\pi^2n)^{1/3} = \qty{1.36e10}{m^{-1}}$~:
$\lambda_{\text{F}} = \qty{4.6}{\angstrom}$. (b) $2a =
\qty{7.2}{\angstrom}$~: le même ordre --- la surface de Fermi du
cuivre s'approche de la frontière de zone, et la physique
d'ouverture du gap se joue \emph{aux} énergies qui comptent. (c) Une
onde stationnaire entasse sa charge sur les ions positifs (énergie
électrostatique plus basse), l'autre entre eux (plus haute)~: même
longueur d'onde, deux énergies --- le gap. (d) Le libre parcours de
110 \omterm{def:b3:crystalline-solids:lattice}{paramètres de maille}~: les ondes de Bloch ne diffusent pas sur
un \omterm{def:b3:crystalline-solids:lattice}{réseau} parfait, il ne reste donc que les vibrations et les
impuretés pour résister.
\end{solution}

\begin{solution}{exo:b3:electrons-in-solids:8}
(a) $E_{\text{g}}/2k_{\text{B}} = \qty{6500}{K}$~:
$\exp[6500(1/300 - 1/400)] \approx 230$. (b) La résistance de la
thermistance plonge exponentiellement --- une courbe raide et
sensible~; celle du platine monte doucement et linéairement (plus de
phonons). L'une est un thermomètre par comptage de porteurs, l'autre
par diffusion. (c) Les porteurs naissent en \emph{paires}
électron--trou, et l'équilibre $np = n_{\text{i}}^2$ partage le coût
du gap entre les deux --- d'où le $E_{\text{g}}/2$. (d)
$n_{\text{i}}$ atteint $\qty{5e21}{m^{-3}}$ vers \qty{760}{K}~: le
\omterm{def:b3:electrons-in-solids:doping}{dopage} est noyé et le circuit oublie sa conception. Les composants
réels abandonnent plus tôt --- leurs fuites inverses, qui doublent
tous les dix kelvins, se conduisent mal les premières.
\end{solution}

\begin{solution}{exo:b3:electrons-in-solids:9}
(a) $\qty{5e22}{m^{-3}}$ --- cinq millions de fois $n_{\text{i}}$.
(b) Le même facteur $\sim5\times10^6$~: une poussière par million
d'atomes s'approprie la conductivité. (c) Parce que des ppm
accidentelles feraient de même sans y être invitées~: seul un
\omterm{def:b3:crystalline-solids:lattice}{cristal} pur au ppb près a une conductivité qui appartient au
concepteur. (d) $d \sim n^{-1/3}
\approx \qty{27}{nm}$, dix fois l'orbite de \qty{2.4}{nm}~: les
\omterm{def:b3:electrons-in-solids:doping}{donneurs} sont des hydrogènes isolés~; pousser le \omterm{def:b3:electrons-in-solids:doping}{dopage} d'un facteur
cent et les orbites se touchent --- une bande d'impuretés, et
finalement un \omterm{def:b3:electrons-in-solids:classes}{métal}.
\end{solution}

\begin{solution}{exo:b3:electrons-in-solids:10}
(a) $E = 13.6\times0.26/11.7^2 \approx \qty{26}{meV}$ (la valeur
mesurée pour le phosphore, \qty{45}{meV}, garde l'échelle honnête).
(b) Comparable à $k_{\text{B}}T = \qty{26}{meV}$~: pratiquement
tous les \omterm{def:b3:electrons-in-solids:doping}{donneurs} sont ionisés à température ambiante --- la
prémisse de \cref{def:b3:electrons-in-solids:doping}. (c) $a =
\qty{0.53}{\angstrom}\times\varepsilon_{\text{r}}/(m^*/
m_{\text{e}}) \approx \qty{2.4}{nm}$. (d) Une orbite qui s'étend sur
des dizaines de \omterm{def:b3:crystalline-solids:lattice}{mailles} voit le \omterm{def:b3:crystalline-solids:lattice}{cristal} comme un milieu lisse ---
et c'est précisément la condition sous laquelle un
$\varepsilon_{\text{r}}$ de volume et un $m^*$ moyenné sur la bande
sont légitimes~: l'approximation se certifie elle-même.
\end{solution}

\begin{solution}{exo:b3:electrons-in-solids:11}
(a) $V_{\text{bi}} = 0.0259\ln(10^{12}) \approx \qty{0.71}{V}$.
(b) $eV/k_{\text{B}}T = 19.3$~: rapport $\eu^{19.3} \approx
2\times10^{8}$ --- la diode est une voie à sens unique à huit
chiffres. (c) $I_0 \propto n_{\text{i}}^2 \propto
\eu^{-E_{\text{g}}/k_{\text{B}}T}$~: la même exponentielle qu'en
\cref{exo:b3:electrons-in-solids:8}, d'où le doublement en règle
empirique. (d) Les quatre diodes forment un losange~: à chaque
demi-période, la paire polarisée en direct dirige le courant dans le
\emph{même} sens à travers la charge --- de l'alternatif à l'entrée,
du continu cahoteux à la sortie.
\end{solution}

\begin{solution}{exo:b3:electrons-in-solids:12}
(a) Rétrécir le gap et davantage de photons le franchissent, mais
chacun ne délivre que la petite énergie de gap~; l'élargir et
chaque photon paie plus cher mais moins nombreux sont admis~: le
produit récolte $\times$ tension culmine entre les deux.
(b) Juste sous l'optimum de \qty{1.3}{eV}~: le plafond idéal du
silicium est de $\sim30\,\%$ --- assez proche pour que son bas prix
l'emporte. (c) La cellule supérieure à grand gap prend le bleu à
haute tension et laisse passer le rouge vers la cellule à petit gap
en dessous~: chaque photon est récolté près de sa propre énergie, la
thermalisation diminue, et le plafond de l'empilement grimpe vers
les quarante. (d) La chaleur relève $I_0$ exponentiellement, ce qui
tire vers le bas la tension en circuit ouvert
$\sim(k_{\text{B}}T/e)\ln(I_{\text{L}}/I_0)$~; et le gap lui-même
rétrécit légèrement, échangeant de la tension contre un courant
qu'il ne peut pas entièrement récupérer.
\end{solution}

\begin{solution}{pb:b3:electrons-in-solids:1}
\textbf{1.} $\lambda_{\text{pic}} \approx \qty{500}{nm}$~:
environ \qty{2.5}{eV}.
\textbf{2.} $hc/E_{\text{g}} \approx \qty{1.1}{\micro m}$.
\textbf{3.} Elle traverse les cellules (aucun état final où
s'absorber) et finit en chaleur dans le support --- en réchauffant
le toit, non les fils.
\textbf{4.} $\qty{1.12}{eV}$ survivent en paire~; l'excédent de
\qty{1.4}{eV} s'écoule en phonons en quelques picosecondes --- plus
vite qu'aucun circuit ne pourrait l'intercepter.
\textbf{5.} La transparence sous le gap ($\sim20\,\%$ de la
puissance) et la thermalisation au-dessus du gap
($\sim30\,\%$)~: le spectre est taxé de moitié avant que la jonction
ne voie le moindre électron.
\textbf{6.} L'absorption est un saut quantique à un photon qui
réclame un état final réel~; à la moitié du gap il n'y en a aucun où
faire halte, et les événements à deux photons sont négligeables aux
densités de photons du soleil.
\textbf{7.} Les paires créées dans la \omterm{prop:b3:electrons-in-solids:junction}{zone de déplétion} ou tout près
ressentent le champ interne~: les électrons sont balayés vers le
côté n, les \omterm{def:b3:electrons-in-solids:doping}{trous} vers le côté p, et le champ les sépare plus vite
qu'ils ne peuvent se retrouver pour se recombiner --- la charge
s'accumule jusqu'à ce qu'un circuit extérieur la soulage en courant.
\textbf{8.} $V_{\text{bi}} = 0.0259\ln(10^{12}) \approx
\qty{0.71}{V}$.
\textbf{9.} $60\times\qty{0.6}{V} = \qty{36}{V}$.
\textbf{10.} $I = 400/36 \approx \qty{11}{A}$.
\textbf{11.} À l'échelle du panneau, la lumière au-dessus du gap
vaut $\sim\qty{1400}{W}$ à $\sim\qty{2}{eV}$ par photon~:
$4\times10^{21}$ photons par seconde --- mais les soixante cellules
sont en \emph{série}, donc les \qty{11}{A} (un flux d'électrons $I/e \approx
7\times10^{19}\,\unit{s^{-1}}$) traversent chaque cellule, dont la
part propre de photons vaut $4\times10^{21}/60 \approx
7\times10^{19}$ par seconde~: presque chaque photon absorbé donne un
électron collecté. Le rendement quantique est excellent~; les pertes
sont énergétiques, non numériques.
\textbf{12.} Après les 50\,\% du spectre, la cellule paie le déficit
de tension ($0.6/1.12 \approx 54\,\%$) puis quelques points de
réflexion et de recombinaison~: $0.5\times0.54 \approx 27\,\%$,
rognés jusqu'aux 20\,\% nominaux.
\textbf{13.} La fuite $I_0$~: la tension en circuit ouvert
$(k_{\text{B}}T/e)\ln(I_{\text{L}}/I_0)$ s'arrête là où le
photocourant et la fuite de la diode s'équilibrent --- environ
\qty{0.6}{V}, bien en deçà du gap.
\textbf{14.} $\sin30^\circ = 0.5$~: la moitié du nominal de midi par
la seule géométrie --- avant les nuages.
\textbf{15.} $\Delta T = \qty{40}{K}$~: $-16\,\%$, ce qui laisse
$\approx\qty{335}{W}$ --- les journées les plus ensoleillées ne sont
pas les meilleures par watt.
\textbf{16.} La chaleur gonfle $I_0$ exponentiellement, et la
tension en circuit ouvert --- reproche d'un logarithme --- perd une
fraction de pour cent par kelvin~; le gap légèrement rétréci achève
le travail.
\textbf{17.} Le câblage en série impose un courant commun unique~:
la cellule ombragée, qui n'en produit aucun, est mise en
polarisation inverse par ses cinquante-neuf collègues et dissipe
leur puissance en point chaud --- la chaîne s'étrangle à la cellule
la plus faible.
\textbf{18.} La diode de dérivation offre au courant un chemin
direct contournant la sous-chaîne ombragée, en sacrifiant sa tension
plutôt que la production de tout le panneau.
\textbf{19.} $\qty{400}{W}\times0.15\times\qty{8760}{h} \approx
\qty{530}{kWh}$ par an.
\textbf{20.} $\approx\qty{130}{euros}$ par an~: un retour en
environ $300/130 \approx 2.5$ ans, puis deux décennies de profit.
\textbf{21.} $500/530 \approx$ un an~: le panneau rembourse son
énergie de fabrication à peu près aussi vite que son prix.
\textbf{22.} La noirceur n'est pas un choix mais une \omterm{thm:b3:electrons-in-solids:bloch}{structure de
bandes}~: l'absorption réclame un état vide à une énergie de photon
au-dessus d'un état plein, et sous le gap le silicium n'en a
simplement aucun à offrir --- la peinture n'ajoute pas d'états.
\textbf{23.} Empiler au-dessus une cellule à grand gap pour prendre
le bleu à haute tension et laisser descendre le rouge~: le tandem
bat le plafond de la jonction unique. L'obstacle~: l'empilement en
série doit accorder les courants (et les prix) entre les couches.
\textbf{24.} Tout ce qui raccourcit $\tau$ et empoisonne les
jonctions~: les défauts créés par l'ultraviolet et les impuretés qui
diffusent vers l'intérieur diffusent et piègent les porteurs,
l'humidité corrode les contacts, et les points chauds vieillissent
les diodes --- la lente entropie d'un \omterm{def:b3:crystalline-solids:lattice}{cristal} qu'on prie de rester
au soleil trente ans.
\textbf{25.} Le gap récolte la moitié du soleil~; la jonction change
les paires en \qty{36}{V} de courant ordonné~; chaînes en série,
ombre et chaleur d'été prélèvent leur part~; et à
$\sim\qty{530}{kWh}$ par an, contre un retour d'un an en énergie et
de trois ans en argent, l'arbitre tranche~: couvrir le toit.
\end{solution}
